% !TITLE 五代十国
% !SUMMARY 乱世中的词心
% !ORDER 0

\begin{intro}
五代十国是一个乱世。唐王朝灭亡之后，中原相继出现五个短命朝廷，南方则有十个割据政权，你方唱罢我登场，平均一个朝代活不过二十年。

但文学的命运常常和历史本身的命运背道而驰。正是在这短短五六十年的动荡之中，一种新的文学样式——词——完成了一次深刻的蜕变。词在唐代已经出现了，盛唐有李白写过《菩萨蛮》和《忆秦娥》，晚唐有温庭筠把词写得浓艳绮丽。但那时候的词毕竟是"诗余"，是文人在正经的诗之外，偶尔为之的小品。到了五代，词才开始获得了独立的文学地位。这主要归功于两个人：一个是李煜，一个是冯延巳。

李煜是南唐后主，他的词可以清晰地分为两个时期。前半生他是帝王，住在金陵的深宫里，写的词也带着帝王才有的富贵气——但这种富贵气里却有一种少有的真诚。他的《玉楼春》写夜宴的奢华，却以一句"待踏马蹄清夜月"收束，那种对美的迷恋不只是感官的享受，更是一种对生活的深情。后半生他成了亡国之君，被押送到汴京，做了赵宋的俘虏。也正是这一巨变，让他的词发生了质变——从此他写的词不再是帝王的悠闲自适，而是亡国奴的血泪悲歌。"问君能有几多愁，恰似一江春水向东流"——这样的句子，没有亡国的经历是写不出来的。

冯延巳是南唐宰相，但他的词与身份无关，只与心有关。"独立小桥风满袖，平林新月人归后"——他写的是每一个普通人都会感受到的惆怅。他的词里没有具体的人、具体的场景，只有一片弥漫的愁绪。这种愁绪不是因某事而起，也不因某事而消，它就是人生的底色。

从李煜到冯延巳，词完成了从描写外部世界到探索内心世界的转变。这一转变，为宋代的晏殊、欧阳修、苏轼们铺好了路。没有五代的铺垫，就没有宋词的辉煌。
\end{intro}
